% Options for packages loaded elsewhere
\PassOptionsToPackage{unicode}{hyperref}
\PassOptionsToPackage{hyphens}{url}
%
\documentclass[
  11pt,
]{article}
\usepackage{lmodern}
\usepackage{amssymb,amsmath}
\usepackage{ifxetex,ifluatex}
\ifnum 0\ifxetex 1\fi\ifluatex 1\fi=0 % if pdftex
  \usepackage[T1]{fontenc}
  \usepackage[utf8]{inputenc}
  \usepackage{textcomp} % provide euro and other symbols
\else % if luatex or xetex
  \usepackage{unicode-math}
  \defaultfontfeatures{Scale=MatchLowercase}
  \defaultfontfeatures[\rmfamily]{Ligatures=TeX,Scale=1}
\fi
% Use upquote if available, for straight quotes in verbatim environments
\IfFileExists{upquote.sty}{\usepackage{upquote}}{}
\IfFileExists{microtype.sty}{% use microtype if available
  \usepackage[]{microtype}
  \UseMicrotypeSet[protrusion]{basicmath} % disable protrusion for tt fonts
}{}
\makeatletter
\@ifundefined{KOMAClassName}{% if non-KOMA class
  \IfFileExists{parskip.sty}{%
    \usepackage{parskip}
  }{% else
    \setlength{\parindent}{0pt}
    \setlength{\parskip}{6pt plus 2pt minus 1pt}}
}{% if KOMA class
  \KOMAoptions{parskip=half}}
\makeatother
\usepackage{xcolor}
\IfFileExists{xurl.sty}{\usepackage{xurl}}{} % add URL line breaks if available
\IfFileExists{bookmark.sty}{\usepackage{bookmark}}{\usepackage{hyperref}}
\hypersetup{
  hidelinks,
  pdfcreator={LaTeX via pandoc}}
\urlstyle{same} % disable monospaced font for URLs
\usepackage[margin=1in]{geometry}
\setlength{\emergencystretch}{3em} % prevent overfull lines
\providecommand{\tightlist}{%
  \setlength{\itemsep}{0pt}\setlength{\parskip}{0pt}}
\setcounter{secnumdepth}{-\maxdimen} % remove section numbering
\usepackage{mathrsfs}
\usepackage{esint}
\setmathfont{Asana Math}
\setlength{\emergencystretch}{3em}

\author{}
\date{23 September 2026}

\begin{document}

\hypertarget{simultaneous-endpoint-saturation-and-the-positive-inertia-of-chern-periods}{%
\section{Simultaneous endpoint saturation and the positive inertia of
Chern
periods}\label{simultaneous-endpoint-saturation-and-the-positive-inertia-of-chern-periods}}

23 September 2026. Baseline: \texttt{8113dd1}. This note deduces a
global constraint on actual Chern classes from the bounded-form endpoint
theorem. The conclusion concerns the positive part of the intersection
form. No equal-square, nonnegative-square, topological-square, or
negative-square budget is imposed on the classes.

\hypertarget{statement-from-actual-geometric-inputs}{%
\subsection{1. Statement from actual geometric
inputs}\label{statement-from-actual-geometric-inputs}}

Let \texttt{(M,g)} be a closed connected oriented smooth four-manifold,
with the Einstein target normalized to \texttt{Ric\ =\ -g}. Let
\texttt{a\_1,...,a\_m}, \texttt{m\ \textgreater{}=\ 2}, be real first
Chern classes of specified Spin\^{}c structures whose unperturbed
Seiberg--Witten equations are solvable for every metric conformal to
\texttt{g}. Each structure is held fixed during conformal variation. The
classes need not be linearly independent or have equal squares. Put

\[
 V=\operatorname{Vol}(g),\quad
 p_i=(a_{i,g}^+)^2>0,\quad
 r=\fint|\operatorname{Ric}+g|^2,\quad w=\sqrt6|W^+|,
\] \[
 \delta_i=\frac32\int(s^2+2w^2)-72\pi^2p_i\ge0,
 \qquad \kappa=\max_i\{r+\delta_i/V\}\le1/2880.
 \tag{1}
\] Let \texttt{lambda\_1\textgreater{}0} be the first positive
scalar-Laplacian eigenvalue and set

\[
 G(\kappa)=13\kappa+4\sqrt{2\kappa},\qquad
 J=\left(300\sqrt\kappa+
                  12\sqrt{G(\kappa)/\lambda_1}\right)^2.
 \tag{2}
\] For \texttt{kappa=0}, set \texttt{J=0}.

\textbf{Theorem (signed periods and positive inertia).} There are
constant signs \texttt{s\_1=1,s\_2,...,s\_m} such that, for every
\texttt{i},

\[
 \boxed{\frac{2\pi^2}{V}
             \bigl((a_i-s_i a_1)_g^+\bigr)^2\le J.}
 \tag{3}
\] Consequently, for every real vector \texttt{t} with
\texttt{sum\_i\ t\_i\ s\_i=0},

\[
 \boxed{\frac{2\pi^2}{V}\left(\sum_i t_i a_i\right)^2
       \le J\left(\sum_{i=2}^m|t_i|\right)^2
       \le(m-1)J\sum_i t_i^2.}
 \tag{4}
\] The first square in (3) is the positive Hodge norm; the square in (4)
is the full, indefinite cup-product square.

The following matrix form records more of the structure. Define the
normalized cup and positive-period matrices

\[
 Q_{ij}=\frac{2\pi^2}{V}a_i\cdot a_j,\qquad
 P_{ij}=\frac{2\pi^2}{V}a_{i,g}^+\cdot a_{j,g}^+,
 \qquad q_1=P_{11}.
\] There are positive semidefinite matrices \texttt{R,N} and a real
vector \texttt{alpha} such that

\[
 \boxed{P=q_1\alpha\alpha^{\mathsf T}+R,
 \quad Q=q_1\alpha\alpha^{\mathsf T}+R-N,
 \quad \operatorname{tr}R\le(m-1)J,}
 \tag{5}
\] \[
 \alpha_1=1,\qquad
 |\alpha_i-s_i|\le\sqrt{J/q_1},\qquad
 \frac12\le q_1\le\frac32.
 \tag{6}
\] The positive-period matrix itself has the sharper entrywise estimate

\[
 \boxed{\max_{i,j}|P_{ij}-q_1s_i s_j|\le\kappa/36+2J.}
 \tag{6a}
\] This uses the exact common-functional relation among the class
lengths, not merely Cauchy--Schwarz applied to the \texttt{L\^{}2}
approximation. In particular, if \texttt{nu\_2(Q)} denotes the second
largest eigenvalue,

\[
 \boxed{\nu_2(Q)\le(m-1)J.}
 \tag{7}
\] No smallness of \texttt{J} is needed for these finite inequalities.
They become restrictive when \texttt{kappa} and
\texttt{G(kappa)/lambda\_1} are small. In particular,
\texttt{lambda\_1\ \textgreater{}\textgreater{}\ sqrt(kappa)} is
sufficient along a sequence; a fixed positive lower bound for
\texttt{lambda\_1} is stronger than necessary.

\hypertarget{two-useful-cup-product-obstructions}{%
\subsection{2. Two useful cup-product
obstructions}\label{two-useful-cup-product-obstructions}}

For two classes \texttt{a,c}, (3) immediately gives

\[
 \boxed{\frac{2\pi^2}{V}
       \bigl(a^2+c^2-2|a\cdot c|\bigr)\le J.}
 \tag{8}
\] Thus positive cup square for both \texttt{a-c} and \texttt{a+c}
obstructs simultaneous saturation when the scalar gap is large relative
to the error. This criterion does not require either class individually
to have positive square. It may be vacuous, which is correctly reported
by a nonpositive left side.

More generally, if \texttt{nu\_2(Q)\textgreater{}=theta\textgreater{}0}
and

\[
 \kappa\le\frac{\theta}{360000(m-1)},
\] then (7) gives the explicit alternative

\[
 \boxed{\lambda_1\le\frac{576(m-1)G(\kappa)}{\theta}.}
 \tag{9}
\] Indeed, (7) implies
\texttt{sqrt(theta/(m-1))\ \textless{}=\ 300\ sqrt(kappa)+12\ sqrt(G/lambda\_1)};
the first summand is at most half the left side. If \texttt{kappa=0}, a
positive \texttt{theta} is impossible. These statements obstruct a
proposed metric and its specified actual classes; they do not rule out
all Einstein metrics on an arbitrary smooth manifold.

For fixed classes and bounded volume, a positive second eigenvalue of
their unnormalized cup matrix therefore excludes a sequence with
\texttt{kappa\ -\textgreater{}\ 0} and
\texttt{G(kappa)/lambda\_1\ -\textgreater{}\ 0}. If volume grows, the
necessary normalization in (7) must be retained: positive inertia alone,
without a quantitative positive eigenvalue divided by volume, is
insufficient.

\hypertarget{an-elementary-phase-lemma-with-a-squared-error}{%
\subsection{3. An elementary phase lemma with a squared
error}\label{an-elementary-phase-lemma-with-a-squared-error}}

All expectations below use normalized volume. Let \texttt{u,v} be smooth
sections of the same Euclidean rank-three bundle, with
\texttt{\textbar{}u\textbar{},\textbar{}v\textbar{}\textless{}=sqrt(6)}.
Write

\[
 a=|u|^2,\quad b=|v|^2,\quad h=\langle u,v\rangle,\quad
 \mathcal Q(u)=u\otimes u-\frac{|u|^2}{3}I,
 \quad \sigma=\operatorname{sign}h\in\{-1,1\},
\] resolving \texttt{h=0} by \texttt{sigma=1}. Pointwise,

\[
 |\mathcal Q(u)-\mathcal Q(v)|^2
       =\frac23(a-b)^2+2(ab-h^2).                 \tag{10}
\] Two useful estimates, including at zeros, are

\[
 |u-\sigma v|^2
 \le2\bigl((a-1)^2+(b-1)^2\bigr)
                    +2|\mathcal Q(u)-\mathcal Q(v)|^2,\tag{11}
\] \[
 |h-\sigma|^2
 \le2\bigl((a-1)^2+(b-1)^2\bigr)
                       +|\mathcal Q(u)-\mathcal Q(v)|^2.\tag{12}
\]

For (11), if \texttt{a+b\textgreater{}=1}, use

\[
 |u-\sigma v|^2
 =\frac{(a-b)^2+4(ab-h^2)}{a+b+2|h|}
 \le2|\mathcal Q(u)-\mathcal Q(v)|^2.
\] If \texttt{a+b\textless{}1}, the left side is at most \texttt{a+b},
whereas
\texttt{2((a-1)\^{}2+(b-1)\^{}2)\textgreater{}=(a+b-2)\^{}2\textgreater{}=a+b}.
For (12), let \texttt{z=sqrt(ab)} and, when \texttt{z\textgreater{}0},
\texttt{t=\textbar{}h\textbar{}/z}. Then
\texttt{0\textless{}=t\textless{}=1} and

\[
 (zt-1)^2\le2(z-1)^2+2z^2(1-t)^2,
 \quad (z-1)^2\le(a-1)^2+(b-1)^2,
 \quad 2z^2(1-t)^2\le2(ab-h^2).
\] The middle inequality follows because \texttt{sqrt(ab)} lies between
\texttt{a} and \texttt{b}; the zero case follows directly. Equation (10)
completes the proof.

Suppose \texttt{E\textbar{}h-sigma\textbar{}\^{}2\textless{}=T} and
\texttt{E\textbar{}dh\textbar{}\^{}2\textless{}=D}. Poincare and the
triangle inequality for orthogonal projection off the constants give

\[
 \sqrt{\operatorname{Var}\sigma}
 \le\sqrt{\operatorname{Var}h}+\|h-\sigma\|_2
 \le\sqrt{D/\lambda_1}+\sqrt T.                  \tag{13}
\] There is a constant sign \texttt{s} with

\[
 \mathbb E|\sigma-s|^2
 =2(1-|\mathbb E\sigma|)
 \le2\operatorname{Var}\sigma.                  \tag{14}
\] If \texttt{E\textbar{}u-sigma\ v\textbar{}\^{}2\textless{}=F}, it
follows that

\[
 \|u-sv\|_2
 \le\sqrt F+\sqrt{12}\bigl(\sqrt{D/\lambda_1}+\sqrt T\bigr).
 \tag{15}
\] All norms here are normalized \texttt{L\^{}2} norms. No derivative of
the measurable sign, unit eigenvector, or normalized form is taken. The
argument only needs the Poincare estimate for these finitely many smooth
scalar functions \texttt{h}; a gap on the whole scalar spectrum is a
convenient sufficient condition.

\hypertarget{applying-the-lemma-to-the-actual-endpoint-forms}{%
\subsection{4. Applying the lemma to the actual endpoint
forms}\label{applying-the-lemma-to-the-actual-endpoint-forms}}

Let \texttt{H\_i} represent \texttt{a\_\{i,g\}\^{}+} harmonically and
put \texttt{xi\_i=-sqrt(2)\ pi\ H\_i}. The actual bounded-form theorem
in \href{line1_composed_application.md}{the endpoint application},
equations (2)--(4), gives smooth \texttt{u\_i=omega\_i/sqrt(2)} with
arbitrarily small simultaneous smoothing allowances. Suppressing those
allowances until the limiting step, it gives

\[
 \begin{split}
 |u_i|&\le\sqrt6,\qquad
 \mathbb E|u_i-\xi_i|^2\le49\kappa/2,\\
 \mathbb E(|u_i|^2-1)^2&\le180\kappa,\qquad
 \mathbb E|W^++\mathcal Q(u_i)|^2\le720\kappa,\\
 \mathbb E|\nabla u_i|^2&\le G(\kappa)/2.
 \end{split}                                                   \tag{16}
\] The same Weyl tensor occurs for every class, so

\[
 \mathbb E|\mathcal Q(u_i)-\mathcal Q(u_1)|^2\le2880\kappa.
 \tag{17}
\] Equations (11)--(12), and the differential estimate
\texttt{\textbar{}d\textless{}u\_i,u\_1\textgreater{}\textbar{}\^{}2\ \textless{}=12(\textbar{}nabla\ u\_i\textbar{}\^{}2+\textbar{}nabla\ u\_1\textbar{}\^{}2)},
give

\[
 F=6480\kappa,\qquad T=3600\kappa,\qquad D=12G(\kappa).
 \tag{18}
\] Thus (15), followed by the two harmonic-alignment errors, gives a
constant sign \texttt{s\_i} with

\[
 \|\xi_i-s_i\xi_1\|_2
 \le\left(7\sqrt2+36\sqrt5+120\sqrt3\right)\sqrt\kappa
                  +12\sqrt{G(\kappa)/\lambda_1}
 \le300\sqrt\kappa+12\sqrt{G(\kappa)/\lambda_1}.       \tag{19}
\] The coefficient estimate is elementary: use
\texttt{sqrt(2)\textless{}3/2}, \texttt{sqrt(5)\textless{}9/4}, and
\texttt{sqrt(3)\textless{}26/15}, giving an upper bound \texttt{299.5}.

To restore the allowances rigorously, choose a sequence of positive
allowances tending to zero, obtaining one finite sign vector for each
choice. Pass to a subsequence with constant sign vector. The harmonic
forms, their norms, and the scalar eigenvalue are fixed. All the errors
in (16)--(19) tend to their displayed bounds. No convergence of spinors,
connections, or measurable sign functions is asserted. This proves (3).
Since a full cup square is at most its positive Hodge square, summing
the signed differences proves (4).

For completeness, (16) implies

\[
 \left(\sqrt{1-\sqrt{180\kappa}}-7\sqrt{\kappa/2}\right)^2
 \le q_1\le
 \left(\sqrt{1+\sqrt{180\kappa}}+7\sqrt{\kappa/2}\right)^2.
 \tag{20}
\] The lower parenthesis is positive on
\texttt{kappa\textless{}=1/2880}, and this interval gives
\texttt{1/2\textless{}=q\_1\textless{}=3/2}.

\hypertarget{the-matrix-decomposition-and-its-precise-limitation}{%
\subsection{5. The matrix decomposition and its precise
limitation}\label{the-matrix-decomposition-and-its-precise-limitation}}

Use the Hilbert space of self-dual forms with normalized \texttt{L\^{}2}
inner product. Orthogonally split

\[
 \xi_i=\alpha_i\xi_1+\eta_i,\qquad
 \alpha_i=\frac{\langle\xi_i,\xi_1\rangle_{L^2(d\mu/V)}}{q_1},
 \quad \langle\eta_i,\xi_1\rangle=0.
\] Here \texttt{alpha\_1=1}, \texttt{eta\_1=0}, and (3) gives
\texttt{\textbar{}\textbar{}eta\_i\textbar{}\textbar{}\_2\^{}2\textless{}=J}
and \texttt{\textbar{}alpha\_i-s\_i\textbar{}\textless{}=sqrt(J/q\_1)}.
Set \texttt{R\_ij=\textless{}eta\_i,eta\_j\textgreater{}} and let
\texttt{N} be the normalized positive Gram matrix of the anti-self-dual
harmonic representatives. Then \texttt{R,N\textgreater{}=0},
\texttt{tr\ R\textless{}=\ (m-1)J}, and Hodge decomposition gives (5)
exactly. On the codimension-one subspace perpendicular to
\texttt{alpha}, the quadratic form of \texttt{Q} is at most
\texttt{\textbar{}\textbar{}R\textbar{}\textbar{}\_op\textless{}=tr\ R}.
The min--max principle proves (7).

To prove (6a), put \texttt{q\_i=P\_ii}. Subtracting the definitions of
the deficits in (1) gives the exact identity

\[
 q_i-q_1=-\frac{\delta_i-\delta_1}{36V},\qquad
 |q_i-q_1|\le\kappa/36.                            \tag{21}
\] The signed forms \texttt{s\_i\ xi\_i} are each within
\texttt{sqrt(J)} of \texttt{xi\_1}, so the squared distance between any
two is at most \texttt{4J}. Polarization therefore gives

\[
 q_1-\kappa/36-2J\le s_i s_jP_{ij}
                       \le q_1+\kappa/36,
\] proving (6a). No extra square root of \texttt{J} is lost in the
positive Gram matrix estimate.

It would be incorrect to delete \texttt{N} or replace (5) by closeness
of \texttt{Q} to a rank-one sign matrix. The bounded-form estimates
control the positive projections. For example, in an abstract
intersection space of signature \texttt{(1,m)}, choose an orthonormal
positive vector \texttt{e\_0} and negative vectors \texttt{e\_i}, and
let \texttt{a\_i=e\_0+L\ e\_i}. Their positive Gram matrix is exactly
\texttt{11\^{}T}, while their full Gram matrix is
\texttt{11\^{}T-L\^{}2\ I}. This illustrates the missing negative
information even when positive-period alignment is exact. It is an
algebraic test of the implication, \textbf{not} an assertion that these
classes are actual monopoles on a metric saturating (1). Classical gauge
equality rigidity can impose further restrictions at exact saturation.

Similarly, distinct Spin\^{}c labels, torsion differences, or the pair
\texttt{a,-a} do not by themselves supply a positive period obstruction.
An integral lattice alone gives no lower bound on the positive Hodge
norm of all nonzero vectors; one must use a verified positive cup
square, a quantitative positive eigenvalue, or additional period data.

\hypertarget{exact-increment-over-the-earlier-arguments}{%
\subsection{6. Exact increment over the earlier
arguments}\label{exact-increment-over-the-earlier-arguments}}

\href{round18_periods.md}{Round 18}, Sections 1--4, assumes common
positive square \texttt{a\_i\^{}2=A\textgreater{}=2chi+3tau} and uses
scalar-energy saturation to control the \textbf{full} determinant
curvatures. It obtains an intersection sign-correlation polytope and low
scalar eigenvalues from covariance. Those remain stronger in their
particular regime and are not new here. \href{round2_rigidity.md}{Round
2}, Sections 3--4, already handles measurable line alignment and Walsh
phase separation at a different calibration.

The added result is (3)--(7) for the actual all-scalar endpoint input,
with arbitrary individual squares and no negative-square budget. It
extracts a positive-inertia obstruction from harmonic Chern periods,
while preserving the uncontrolled negative Gram matrix. The squared
phase estimates (11)--(14) yield an \texttt{O(kappa+G(kappa)/lambda\_1)}
error, rather than losing a square root simply by estimating the mean of
the inner product. The common-functional identity (21) preserves that
order in the positive Gram matrix. No growing low-mode cluster theorem
is being claimed.

The classical scalar--Weyl curvature bound and its equality geometry
come from C. LeBrun, \emph{Ricci Curvature, Minimal Volumes, and
Seiberg--Witten Theory},
\href{https://arxiv.org/abs/math/0003068}{Theorem 2.3 and Proposition
4.1}. The local endpoint application provides the quantitative
all-scalar harmonic alignment used here.

\texttt{line1\_period\_inertia\_check.py} checks the exact tensor
identities, phase inequalities' algebra, constants, and a finite model
of (5). It does not certify gauge solvability or replace Poincare, Hodge
theory, the smoothing argument, or the analytic endpoint theorem.

\end{document}
